A degraded register is a state of the world and not a training accident, so
\textbf{the mechanism is applied to both halves of the data}: the same values
are missing, stale, corrupted or split at fit time and at evaluation time, and
the increment measures a model built on a degraded register and deployed
against one. What is estimated from the \emph{training half alone} is the
mechanism's \emph{parameters} --- which values fall in the long tail, which are
common, the frequency table a corruption draws from, the set of identities a
reconciliation failure splits --- so that no test outcome and no test-half
content can reach the counterfactual.

That property is executed rather than asserted: the test half is perturbed,
the mechanism is re-applied, and the training half's degraded values are
required to be identical, on every mechanism and level. The check runs on
every mechanism and every level in the design space and is one of the
conditions \texttt{verify\_numbers} enforces
(Supplement~\ref{app:verifier}).

The distinction matters because the two designs measure different things. A
mechanism applied to the training half alone would degrade the model's inputs
and leave the evaluation clean, which measures a train--test mismatch: how
much a model built on a bad register loses when deployed against a good one.
That is not the question a register's owner is asking. The owner's register is
bad at both moments, and the counterfactual that answers the question degrades
both.
